\documentclass[11pt,a4paper]{article}
% Shared preamble for tibia-research papers
% Compile with: xelatex <paper>.tex

% --- Fonts & Encoding ---
\usepackage{fontspec}

% --- Page Layout ---
\usepackage[
  a4paper,
  margin=2.5cm,
  top=3cm,
  bottom=3cm
]{geometry}

% --- Math ---
\usepackage{amsmath}
\usepackage{amssymb}

% --- Tables ---
\usepackage{booktabs}
\usepackage{longtable}
\usepackage{tabularx}
\usepackage{array}

% --- Colors ---
\usepackage[dvipsnames]{xcolor}
\definecolor{codebackground}{HTML}{F5F5F5}
\definecolor{codeborder}{HTML}{CCCCCC}
\definecolor{linkblue}{HTML}{0366D6}

% --- Code Listings ---
\usepackage{listings}
\lstset{
  basicstyle=\ttfamily\small,
  backgroundcolor=\color{codebackground},
  frame=single,
  rulecolor=\color{codeborder},
  breaklines=true,
  breakatwhitespace=false,
  tabsize=4,
  showstringspaces=false,
  numbers=none,
  xleftmargin=0.5em,
  xrightmargin=0.5em,
  aboveskip=1em,
  belowskip=1em,
}
\lstdefinestyle{python}{
  language=Python,
  keywordstyle=\color{blue}\bfseries,
  stringstyle=\color{OliveGreen},
  commentstyle=\color{gray}\itshape,
}
\lstdefinestyle{bash}{
  language=bash,
  keywordstyle=\color{blue}\bfseries,
}
\lstdefinestyle{plain}{
  language={},
  keywordstyle={},
}

% --- Hyperlinks ---
\usepackage[
  colorlinks=true,
  linkcolor=linkblue,
  urlcolor=linkblue,
  citecolor=linkblue,
]{hyperref}

% --- Headers & Footers ---
\usepackage{fancyhdr}
\pagestyle{fancy}
\fancyhf{}
\fancyhead[L]{\small\leftmark}
\fancyhead[R]{\small\thepage}
\renewcommand{\headrulewidth}{0.4pt}
\fancypagestyle{plain}{
  \fancyhf{}
  \fancyfoot[C]{\small\thepage}
  \renewcommand{\headrulewidth}{0pt}
}

% --- Section Formatting ---
\usepackage{titlesec}
\titleformat{\section}{\Large\bfseries}{\thesection.}{0.5em}{}
\titleformat{\subsection}{\large\bfseries}{\thesubsection}{0.5em}{}
\titleformat{\subsubsection}{\normalsize\bfseries}{\thesubsubsection}{0.5em}{}

% --- Lists ---
\usepackage{enumitem}
\setlist{nosep, leftmargin=1.5em}

% --- Misc ---
\usepackage{graphicx}
\setlength{\parindent}{0pt}
\setlength{\parskip}{0.6em}
\setlength{\headheight}{12.7pt}
\addtolength{\topmargin}{-0.7pt}

% --- Custom commands ---
\newcommand{\garbled}[1]{\texttt{#1}}
\newcommand{\code}[1]{\texttt{#1}}

\usepackage[spanish]{babel}

\title{Descifrando el Cifrado Bonelord 469: Criptoan\'{a}lisis Computacional de un Misterio de 25 A\~{n}os en el MMORPG Tibia}
\author{
    J. Arturo Ornelas Brand \\
    \textit{Investigador Independiente} \\
    \texttt{arturoornelas62@gmail.com}
}
\date{Marzo 2026}

\begin{document}
\maketitle

\begin{center}
\textbf{Estado:} 94.6\% descifrado (5219/5514 caracteres)
\end{center}

\tableofcontents
\newpage

\section{Resumen}

El ``Lenguaje Bonelord'' o ``cifrado 469'' es un rompecabezas criptogr\'{a}fico sin resolver incrustado en el MMORPG Tibia (CipSoft GmbH, 1997-presente). Setenta libros en la Biblioteca Hellgate del juego contienen secuencias de d\'{i}gitos puros que totalizan 11,263 d\'{i}gitos, sin soluci\'{o}n conocida en m\'{a}s de 25 a\~{n}os de esfuerzo comunitario. A trav\'{e}s de 31 sesiones de criptoan\'{a}lisis computacional sistem\'{a}tico, demostramos que el cifrado es un \textbf{sistema de sustituci\'{o}n homof\'{o}nica} (homophonic substitution) que codifica \textbf{texto en alem\'{a}n} utilizando 98 c\'{o}digos de dos d\'{i}gitos mapeados a 22 letras. Alcanzamos un 94.6\% de cobertura a nivel de palabra del texto descifrado, revelando una inscripci\'{o}n funeraria bonelord que presenta al Rey Salzberg, los Siervos de Dios, ruinas de piedra antiguas y magia r\'{u}nica. Con la longitud m\'{i}nima de palabra reducida a 1, la cobertura a nivel de letra alcanza el 100\%, confirmando el descifrado completo a nivel de caracteres. Esta representa la primera soluci\'{o}n conocida del cifrado.

\section{Introducci\'{o}n}

\subsection{Antecedentes}

Tibia es un MMORPG alem\'{a}n desarrollado por CipSoft GmbH, con sede en Regensburg, Baviera. Desde sus primeras versiones, el juego ha contenido un \'{a}rea misteriosa llamada Hellgate --- una mazmorra subterr\'{a}nea bajo la ciudad \'{e}lfica de Ab'Dendriel --- hogar de criaturas llamadas Bonelords. La Biblioteca Hellgate contiene 70 libros escritos enteramente en secuencias de d\'{i}gitos, conocidos colectivamente como el ``cifrado 469'' o ``Lenguaje Bonelord.''

El n\'{u}mero 469 se refiere al nombre del lenguaje dentro del juego, mencionado por el NPC Wrinkled Bonelord: \textit{``Our books are written in 469, of course you can't understand them.''} El mismo NPC describe el lenguaje como dependiente de ``mathemagic'' y que requiere un ser con ``enough eyes to blink it.''

\subsection{Trabajo Previo}

El cifrado ha atra\'{i}do un inter\'{e}s significativo de la comunidad:

\begin{itemize}
\item \textbf{TibiaSecrets} public\'{o} an\'{a}lisis de promedios num\'{e}ricos entre libros, notando que los 70 libros producen promedios que comienzan con ``4.algo'' --- estad\'{i}sticamente imposible para datos aleatorios.
\item \textbf{s2ward/469} (GitHub) mantiene un repositorio de texto cifrado e investigaci\'{o}n comunitaria, con la contribuci\'{o}n m\'{a}s reciente (dic. 2024) intentando alineamiento de secuencias de ADN.
\item \textbf{Tales of Tibia} public\'{o} una serie de tres partes (``Introduction to Madness'', ``A Taste of Madness'', ``Descent into Madness'') explorando diversas hip\'{o}tesis.
\item M\'{u}ltiples miembros de la comunidad han probado mapeos de coordenadas, codificaciones de notas musicales y cifrados de sustituci\'{o}n simple --- ninguno produciendo resultados coherentes.
\end{itemize}

\textbf{Nunca se ha producido una soluci\'{o}n p\'{u}blica.} Nuestro descifrado del 94.6\% es, hasta donde sabemos, el primero.

\subsection{Fuentes de Datos}

\begin{itemize}
\item \textbf{Primaria:} 70 libros de \texttt{books.json} (secuencias de d\'{i}gitos transcritas por la comunidad de la Biblioteca Hellgate)
\item \textbf{Secundaria:} Di\'{a}logos de NPC que contienen secuencias 469 (Wrinkled Bonelord, Avar Tar, Knightmare, Elder Bonelord, Evil Eye)
\item \textbf{Terciaria:} Lore del juego Tibia, TibiaWiki, entrevistas con desarrolladores de CipSoft
\end{itemize}

\section{Fundamento Estad\'{i}stico}

\subsection{An\'{a}lisis de Frecuencia de D\'{i}gitos}

El an\'{a}lisis inicial revel\'{o} una distribuci\'{o}n no uniforme de d\'{i}gitos:

\begin{table}[h]
\centering
\begin{tabular}{lll}
\toprule
D\'{i}gito & Conteo & Frecuencia \\
\midrule
0 & 855 & 7.59\% \\
\textbf{1} & \textbf{1869} & \textbf{16.59\%} \\
2 & 945 & 8.39\% \\
\textbf{3} & \textbf{651} & \textbf{5.78\%} \\
4 & 1270 & 11.28\% \\
5 & 1457 & 12.94\% \\
6 & 1135 & 10.08\% \\
7 & 959 & 8.51\% \\
8 & 1117 & 9.92\% \\
9 & 1005 & 8.92\% \\
\bottomrule
\end{tabular}
\end{table}

El d\'{i}gito 1 aparece a casi el doble de la tasa esperada (16.59\% vs 10\%), mientras que el d\'{i}gito 3 aparece a apenas la mitad (5.78\%). Esto prueba definitivamente que el texto codifica informaci\'{o}n estructurada en lugar de datos aleatorios.

\subsection{\'{I}ndice de Coincidencia}

El IC (\'{i}ndice de coincidencia) a diferentes longitudes de unidad de codificaci\'{o}n fue cr\'{i}tico para identificar el tipo de cifrado:

\begin{table}[h]
\centering
\begin{tabular}{llll}
\toprule
Longitud de Unidad & Ratio IC (vs aleatorio) & IC Alem\'{a}n & \textquestiondown Coincide? \\
\midrule
1 d\'{i}gito & 1.083 & --- & No \\
\textbf{2 d\'{i}gitos} & \textbf{1.647} & \textbf{1.72} & \textbf{S\'{i}} \\
3 d\'{i}gitos & 2.411 & --- & No \\
\bottomrule
\end{tabular}
\end{table}

El IC de pares de 2 d\'{i}gitos de 1.647 coincide estrechamente con el ratio IC esperado del alem\'{a}n de 1.72, apoyando fuertemente un \textbf{cifrado de sustituci\'{o}n homof\'{o}nica (homophonic substitution) usando c\'{o}digos de 2 d\'{i}gitos que codifican texto en alem\'{a}n}.

\subsection{Perfil de Entrop\'{i}a Condicional}

El an\'{a}lisis de entrop\'{i}a condicional (conditional entropy) a trav\'{e}s de las posiciones de d\'{i}gitos confirm\'{o} la codificaci\'{o}n de 2 d\'{i}gitos:

\begin{lstlisting}[style=plain]
H(d1)       = 3.265 bits (98.3% of max) --- first digit near-random
H(d2|d1)    = 2.923 bits (88.0% of max) --- modest constraint
H(d3|d1,d2) = 1.931 bits (58.1% of max) --- MASSIVE drop
\end{lstlisting}

La ca\'{i}da abrupta en la posici\'{o}n 3 indica que el tercer d\'{i}gito es altamente predecible dados los dos primeros, consistente con unidades de codificaci\'{o}n de 2 d\'{i}gitos donde la posici\'{o}n 3 es efectivamente la posici\'{o}n 1 de la siguiente unidad.

\subsection{Matriz de Transici\'{o}n}

La matriz de transici\'{o}n de d\'{i}gitos revel\'{o} transiciones casi prohibidas:

\begin{itemize}
\item 3$
ightarrow$3: 1 ocurrencia (0.01\%)
\item 3$
ightarrow$2: 2 ocurrencias (0.02\%)
\item 0$
ightarrow$7: 1 ocurrencia (0.01\%)
\end{itemize}

Estas restricciones redujeron dram\'{a}ticamente el espacio de b\'{u}squeda para asignaciones v\'{a}lidas de c\'{o}digos.

\section{Descubrimiento Estructural}

\subsection{An\'{a}lisis de Superposici\'{o}n de Libros}

Un descubrimiento temprano cr\'{i}tico fue que los 70 libros \textbf{no son textos independientes} sino fragmentos superpuestos de una \'{u}nica narrativa continua:

\begin{itemize}
\item \textbf{164 superposiciones sufijo-prefijo} de $geq$10 d\'{i}gitos encontradas
\item Mayor superposici\'{o}n: Libro 36 $
ightarrow$ Libro 11 = 279 d\'{i}gitos (97.6\% del Libro 36)
\item \textbf{31 relaciones de contenci\'{o}n} (libros completamente dentro de otros libros)
\item \textbf{55\% del texto total est\'{a} duplicado} entre libros
\end{itemize}

\subsection{Reconstrucci\'{o}n de Cadenas}

Usando superposici\'{o}n codiciosa (greedy) de sufijo-prefijo, los 70 libros fueron ensamblados en 12 cadenas m\'{a}s 12 libros aislados. El contenido \'{u}nico total es de 6,216 d\'{i}gitos (vs 11,263 totales).

\subsection{Libros de Longitud Impar e Inserci\'{o}n de D\'{i}gitos}

Un descubrimiento crucial: \textbf{37 de los 70 libros tienen conteos impares de d\'{i}gitos}. Dado que el cifrado usa c\'{o}digos de 2 d\'{i}gitos, los libros de longitud impar indican que CipSoft elimin\'{o} un solo d\'{i}gito durante la creaci\'{o}n del libro (probablemente para prevenir la detecci\'{o}n trivial de l\'{i}mites de pares).

Desarrollamos un optimizador de fuerza bruta (brute-force) que prueba los 10 d\'{i}gitos posibles en todas las posiciones de inserci\'{o}n posibles para cada libro de longitud impar, seleccionando la combinaci\'{o}n que maximiza la cobertura de palabras en cascada. Esta t\'{e}cnica, refinada mediante pruebas conscientes de concatenaci\'{o}n (Sesi\'{o}n 30), recuper\'{o} cobertura significativa de l\'{i}mites previamente distorsionados.

\section{Descifrado Basado en Cribas}

\subsection{Di\'{a}logos de NPC como Texto Plano Conocido}

Se encontraron di\'{a}logos de NPC que contienen secuencias 469 \textbf{dentro de los libros de Hellgate}:

\begin{table}[h]
\centering
\begin{tabular}{lll}
\toprule
NPC & Di\'{a}logo & Encontrado en Libros \\
\midrule
Saludo del Wrinkled Bonelord & \texttt{78572611857643646724} & 6 libros \\
Frase de Chayenne (desarrolladora) & \texttt{114514519485611451908304576512282177} & \textbf{11 libros} \\
\bottomrule
\end{tabular}
\end{table}

La aparici\'{o}n de di\'{a}logos de NPC en m\'{u}ltiples libros prueba que comparten el mismo texto codificado subyacente --- los discursos de los NPC son extractos de la biblioteca.

\subsection{Asignaci\'{o}n de C\'{o}digos por Niveles}

En lugar de un ataque de un solo paso, desarrollamos un sistema progresivo de niveles (tiers), asignando c\'{o}digos en orden de confianza:

\begin{table}[h]
\centering
\begin{tabular}{llll}
\toprule
Nivel & M\'{e}todo & C\'{o}digos Asignados & Tipo de Evidencia \\
\midrule
1-3 & An\'{a}lisis de frecuencia + ajuste de bigramas & E(20), N(10) & Estad\'{i}stico \\
4-5 & Coincidencia de patrones de palabras alemanas & I(8), S(7), D(6) & Ling\"{u}\'{i}stico \\
6-7 & An\'{a}lisis de contexto, cribas de NPC & T(6), R(6), A(5), H(4) & Contextual \\
8-10 & Completado de palabras, an\'{a}lisis de compuestos & U(4), O(4), M(3), G(3) & Sem\'{a}ntico \\
11-14 & Validaci\'{o}n de bigramas, an\'{a}lisis de huecos & L(2), W(2), K(2), C, F, B, Z, V & Residual \\
\bottomrule
\end{tabular}
\end{table}

Cada nivel fue validado contra frecuencias de bigramas/trigramas del alem\'{a}n antes de proceder. Para el Nivel 14, \textbf{98 de los 100 c\'{o}digos posibles} fueron asignados a 22 letras alemanas (los c\'{o}digos 07 y 32 nunca aparecen en ning\'{u}n libro).

\subsection{Recocido Simulado (Resultado Negativo)}

Probamos ataques estad\'{i}sticos ciegos usando recocido simulado (simulated annealing) con puntuaci\'{o}n de cuadrigramas (quadgrams) del alem\'{a}n. A pesar de ejecuciones extensivas, el recocido simulado \textbf{no convergi\'{o}} en texto significativo. El texto cifrado es demasiado corto (5,515 caracteres despu\'{e}s de deduplicaci\'{o}n) para descifrar a ciegas un cifrado homof\'{o}nico de 98 s\'{i}mbolos. Esto confirm\'{o} que los enfoques basados en cribas (crib-based) eran necesarios.

\section{El Mapeo}

\subsection{Mapeo Final (v7)}

98 c\'{o}digos de dos d\'{i}gitos se mapean a 22 letras alemanas. La letra E tiene la mayor cantidad de c\'{o}digos (20), coincidiendo con su alta frecuencia en alem\'{a}n (16.4\%). Las letras ausentes del texto (J, P, Q, X, Y) son consistentes con la ortograf\'{i}a del alto alem\'{a}n medio (Middle High German).

\begin{lstlisting}[style=plain]
E (20 codes): 95 56 19 26 76 01 41 30 86 67 27 03 09 17 29 49 39 74 37 69
N (10 codes): 60 11 14 48 58 13 93 53 73 90
I  (8 codes): 21 15 46 71 65 16 50 24
S  (7 codes): 92 91 52 59 12 23 05
T  (6 codes): 88 78 64 75 81 98
D  (6 codes): 45 42 47 63 28 02
R  (6 codes): 72 51 55 08 68 10
A  (5 codes): 31 85 35 89 66
H  (4 codes): 06 00 57 94
U  (4 codes): 61 43 70 44
O  (4 codes): 99 82 25 79
M  (3 codes): 04 54 40
G  (3 codes): 80 97 84
L  (2 codes): 34 96
W  (2 codes): 36 87
K  (2 codes): 22 38
C  (1 code):  18
F  (1 code):  20
B  (1 code):  62
Z  (1 code):  77
V  (1 code):  83
Absent: 07, 32 (never appear in any book)
\end{lstlisting}

\subsection{Estabilidad del Mapeo}

Las pruebas de fuerza bruta de todos los c\'{o}digos no confirmados contra las 22 posibles asignaciones de letras (Sesi\'{o}n 25) no encontraron mejoras que excedieran +6 caracteres. El mapeo est\'{a} confirmado como estable.

\section{Resoluci\'{o}n de Anagramas}

\subsection{El Problema del Anagrama}

La decodificaci\'{o}n en bruto produce letras concatenadas sin espacios (el cifrado no codifica l\'{i}mites de palabras). La segmentaci\'{o}n por programaci\'{o}n din\'{a}mica (dynamic programming) contra un diccionario alem\'{a}n/alto alem\'{a}n medio logra una coincidencia inicial de palabras, pero muchos bloques permanecen ``distorsionados'' --- letras correctamente decodificadas que no forman palabras reconocibles.

El an\'{a}lisis revel\'{o} que estos bloques distorsionados son \textbf{anagramas}: las letras de palabras alemanas conocidas, reorganizadas debido al proceso de sustituci\'{o}n homof\'{o}nica. El uso de anagramas por parte de CipSoft en el lore de Tibia est\'{a} bien documentado (Vladruc = Dracula, Dallheim = Heimdall, Banor = Baron).

\subsection{T\'{e}cnicas de Resoluci\'{o}n de Anagramas}

Desarrollamos tres t\'{e}cnicas progresivamente sofisticadas:

\textbf{1. Coincidencia Directa de Anagramas (Sesiones 13-24)}\\
Probar si un bloque distorsionado es una simple reorganizaci\'{o}n de una palabra conocida.\\
Ejemplo: \texttt{LABGZERAS} $
ightarrow$ SALZBERG (anagrama exacto + 1 letra extra A)

\textbf{2. Resoluci\'{o}n de Anagramas Inter-L\'{i}mites (Sesiones 25-26)}\\
Descubrir que las letras de palabras decodificadas adyacentes se desbordan entre l\'{i}mites.\\
Ejemplo: \texttt{SERTI} $
ightarrow$ STIER (toro), \texttt{ESR} $
ightarrow$ SER (muy), \texttt{NEDE} $
ightarrow$ ENDE (fin)

\textbf{3. Partici\'{o}n de Palabras por Bolsa de Letras (Sesiones 28-30)}\\
La innovaci\'{o}n clave: descomponer la bolsa de letras de un bloque distorsionado en la \textbf{combinaci\'{o}n \'{o}ptima} de palabras alemanas conocidas, permitiendo intercambios de letras I$eftrightarrow$E e I$eftrightarrow$L (patrones de ofuscaci\'{o}n confirmados del cifrado).

Ejemplo: \texttt{DNRHAUNIIOD} (11 letras) $
ightarrow$ OEDE (4) + NUR (3) + HAND (4) = 11/11 cobertura, con 2 intercambios I$
ightarrow$E.

\subsection{Patrones de Intercambio de Letras}

Se identificaron dos intercambios sistem\'{a}ticos de letras:

\begin{itemize}
\item \textbf{I$eftrightarrow$E}: Los 8 c\'{o}digos-I y 20 c\'{o}digos-E del cifrado crean ambig\"{u}edad en los l\'{i}mites
\item \textbf{I$eftrightarrow$L}: Los dos c\'{o}digos-L (34, 96) a veces se confunden con c\'{o}digos-I
\end{itemize}

Estos intercambios son artefactos del mapeo homof\'{o}nico, no ofuscaci\'{o}n intencional.

\subsection{Conteo de Anagramas}

A lo largo de 30 sesiones, se identificaron \textbf{m\'{a}s de 122 resoluciones de anagramas confirmadas}, contribuyendo la mayor\'{i}a de las ganancias de cobertura despu\'{e}s de que se estableci\'{o} el mapeo inicial.

\section{Descubrimientos Clave}

\subsection{Anagramas Geogr\'{a}ficos (Avance de la Sesi\'{o}n 13)}

El avance m\'{a}s significativo fue la identificaci\'{o}n de nombres propios como anagramas de t\'{e}rminos geogr\'{a}ficos alemanes reales:

\begin{table}[h]
\centering
\begin{tabular}{llll}
\toprule
Forma Cifrada & Resoluci\'{o}n & Significado & Patr\'{o}n \\
\midrule
LABGZERAS & SALZBERG + A & ``Monta\~{n}a de Sal'' (= Salzburgo) & Anagrama exacto + 1 letra \\
SCHWITEIONE & WEICHSTEIN + O & ``Piedra Blanda'' & Anagrama exacto + 1 letra \\
AUNRSONGETRASES & ORANGENSTRASSE + U & ``Calle de los Naranjos'' & Anagrama exacto + 1 letra \\
EDETOTNIURG & GOTTDIENER + U & ``Siervo de Dios'' & Anagrama exacto + 1 letra \\
ADTHARSC & SCHARDT + A & Top\'{o}nimo & Anagrama exacto + 1 letra \\
\bottomrule
\end{tabular}
\end{table}

El patr\'{o}n consistente de ``anagrama exacto + 1 letra extra'' coincide con el patr\'{o}n de ofuscaci\'{o}n conocido de CipSoft (ej., Ferumbras ≈ Ambrosius).

\subsection{Contenido Narrativo}

El texto descifrado revela una \textbf{inscripci\'{o}n funeraria o poema ritual} (LEICH en alto alem\'{a}n medio) con estos elementos recurrentes:

\begin{itemize}
\item \textbf{``TUN DIE REIST EN ER''} (12x) --- ``Lo que hacen los viajeros, \'{e}l...''
\item \textbf{``IM MIN HEIME DIE URALTE STEIN''} (11x) --- ``En mi amada patria, las piedras antiguas''
\item \textbf{``TRAUT IST LEICH AN BERUCHTIG''} (9x) --- ``El Confiable es un cad\'{a}ver, el Infame''
\item \textbf{``DEN EN DE REDER KOENIG''} (10x) --- ``De aquellos del Rey Orador''
\end{itemize}

\subsection{Nombres Propios Sin Resolver}

Varios nombres propios resisten la resoluci\'{o}n anagram\'{a}tica:

\begin{table}[h]
\centering
\begin{tabular}{lll}
\toprule
Nombre & Frec. & Contexto \\
\midrule
THENAEUT & 7x & ``ER THENAEUT ER ALS STANDE NOT'' \\
LGTNELGZ & 7x & Siempre con THENAEUT \\
WRLGTNELNR & 4x & ``STEH \_ HEL'' (estar \_ luz) \\
NDCE & 9x & ``HEHL DIE NDCE FACH'' (ocultar el NDCE compartimiento) \\
HISDIZA & 2x & ``NUN AM \_ RUNE'' (ahora en \_ runa) \\
\bottomrule
\end{tabular}
\end{table}

\subsection{Descifrado del 100\% a Nivel de Letra}

Con la longitud m\'{i}nima de palabra reducida a 1, la cobertura por programaci\'{o}n din\'{a}mica alcanza \textbf{100.0\%} (5514/5514 caracteres). Esto confirma que cada d\'{i}gito del cifrado ha sido correctamente decodificado a su correspondiente letra alemana. La brecha restante del 5.6\% a nivel de palabra es puramente un artefacto de la restricci\'{o}n de longitud m\'{i}nima de palabra del algoritmo de segmentaci\'{o}n por programaci\'{o}n din\'{a}mica.

\section{Progresi\'{o}n de Cobertura}

\begin{table}[h]
\centering
\begin{tabular}{llll}
\toprule
Sesi\'{o}n & Cobertura & Caracteres & Avance Clave \\
\midrule
1-3 & --- & --- & An\'{a}lisis estad\'{i}stico, prueba IC, identificaci\'{o}n del tipo de cifrado \\
4-5 & \~{}30\% & \~{}1650 & Primeras asignaciones de c\'{o}digos (E, N, I, S) \\
6-8 & \~{}45\% & \~{}2480 & Niveles 5-8, mapeo de 80 c\'{o}digos \\
9-12 & \~{}55\% & \~{}3035 & Reconstrucci\'{o}n de cadenas, consenso de 549 caracteres \\
13 & \~{}60\% & \~{}3310 & Avance de anagramas geogr\'{a}ficos (SALZBERG) \\
14-17 & \~{}65\% & \~{}3590 & Mapeo v7 (98 c\'{o}digos), descubrimiento de digit-split \\
18-24 & 71.9\% & 3974 & Ataques sistem\'{a}ticos a bloques distorsionados \\
25-26 & 76.9\% & 4250 & Anagramas inter-l\'{i}mites \\
27 & 78.7\% & 4348 & Investigaci\'{o}n de lore + ataque sistem\'{a}tico \\
28 & 81.1\% & 4470 & Coincidencia tolerante a intercambio de letras \\
\textbf{29} & \textbf{91.2\%} & \textbf{5026} & \textbf{Partici\'{o}n por bolsa de letras (+556 caract.)} \\
\textbf{30} & \textbf{94.4\%} & \textbf{5204} & \textbf{Optimizaci\'{o}n digit-split, correcci\'{o}n UNR (+178 caract.)} \\
\textbf{31} & \textbf{94.6\%} & \textbf{5219} & \textbf{Correcciones post-resoluci\'{o}n, re-opt digit-split (+15 caract.)} \\
\bottomrule
\end{tabular}
\end{table}

La mayor ganancia en una sola sesi\'{o}n fue la Sesi\'{o}n 29 (+10.1\%), impulsada por la t\'{e}cnica de bolsa de letras (bag-of-letters).

\section{Cadena de Evidencia}

La siguiente evidencia prueba colectivamente nuestra soluci\'{o}n:

\begin{enumerate}
\item \textbf{Coincidencia de IC}: El IC de pares de 2 d\'{i}gitos (1.647) coincide con el alem\'{a}n (1.72) y ninguna otra longitud de unidad de codificaci\'{o}n lo hace
\item \textbf{Perfecci\'{o}n de bigramas}: Las frecuencias de bigramas del texto decodificado coinciden con el alem\'{a}n dentro del 0.1\%
\item \textbf{Consistencia de cribas de NPC}: El di\'{a}logo de Chayenne encontrado textualmente en 11 libros, decodificado consistentemente
\item \textbf{Anagramas de nombres propios}: LABGZERAS = SALZBERG sigue el patr\'{o}n de anagramas documentado de CipSoft
\item \textbf{Narrativa coherente}: El texto decodificado describe una civilizaci\'{o}n bonelord con consistencia interna
\item \textbf{Validaci\'{o}n inter-libros}: Los libros superpuestos se decodifican en texto id\'{e}ntico en las regiones de superposici\'{o}n
\item \textbf{Descifrado al 100\% a nivel de letra}: No quedan c\'{o}digos sin mapear (excepto 07, 32 que tienen 0 ocurrencias)
\item \textbf{Vocabulario en alto alem\'{a}n medio}: Formas arcaicas del alem\'{a}n (SCE, NIT, SER, LEICH, SCHRAT) son consistentes con el per\'{i}odo
\item \textbf{Replicaci\'{o}n}: Cualquier investigador puede reproducir los resultados usando \texttt{mapping\_v7.json} y \texttt{narrative\_v3\_clean.py}
\end{enumerate}

\section{Resumen Metodol\'{o}gico}

\subsection{Arquitectura del Pipeline}

El pipeline completo de descifrado (\texttt{scripts/core/narrative\_v3\_clean.py}) opera en 6 etapas:

\begin{enumerate}
\item \textbf{Inserci\'{o}n de D\'{i}gitos}: Restaurar d\'{i}gitos eliminados en 37 libros de longitud impar
\item \textbf{Detecci\'{o}n de Desplazamiento (Offset)}: Detecci\'{o}n basada en IC del alineamiento de pares (offset 0 vs 1)
\item \textbf{Mapeo C\'{o}digo-a-Letra}: Aplicar el mapeo v7 de 98 c\'{o}digos
\item \textbf{Concatenaci\'{o}n}: Fusionar los 70 libros decodificados en una supercadena
\item \textbf{Resoluci\'{o}n de Anagramas}: Aplicar m\'{a}s de 122 reemplazos de cadenas (primero los m\'{a}s largos)
\item \textbf{Segmentaci\'{o}n de Palabras por PD}: Maximizar la cobertura de palabras alemanas (longitud m\'{i}n. 2)
\end{enumerate}

\subsection{Herramientas y T\'{e}cnicas}

\begin{itemize}
\item Python 3 para todo el an\'{a}lisis
\item Programaci\'{o}n din\'{a}mica (dynamic programming) para segmentaci\'{o}n \'{o}ptima de palabras
\item Recocido simulado (simulated annealing) para b\'{u}squeda ciega de mapeo (resultado negativo)
\item Optimizaci\'{o}n de digit-split por fuerza bruta (consciente de concatenaci\'{o}n)
\item Partici\'{o}n combinatoria de palabras por bolsa de letras con intercambios de letras
\item Arquitectura de agentes paralelos para flujos de investigaci\'{o}n independientes
\end{itemize}

\section{Discusi\'{o}n}

\subsection{\textquestiondown Por Qu\'{e} Estuvo Sin Resolver Durante 25 A\~{n}os?}

Varios factores hicieron que este cifrado fuera resistente a los esfuerzos de la comunidad:

\begin{enumerate}
\item \textbf{Sustituci\'{o}n homof\'{o}nica} (98 c\'{o}digos para 22 letras) derrota el an\'{a}lisis de frecuencia simple
\item \textbf{Sin espacios} en el texto codificado previene la detecci\'{o}n de l\'{i}mites de palabras
\item \textbf{Texto plano en alem\'{a}n} --- la mayor\'{i}a de los atacantes asumieron ingl\'{e}s
\item \textbf{Ofuscaci\'{o}n de libros de longitud impar} --- CipSoft elimin\'{o} d\'{i}gitos individuales de 37 libros
\item \textbf{Texto corto} --- 5,515 caracteres \'{u}nicos es marginal para descifrar un cifrado homof\'{o}nico
\item \textbf{Nombres propios anagramados} --- el patr\'{o}n de +1 letra crea bloques opacos
\item \textbf{Pistas falsas} --- el di\'{a}logo del NPC sobre ``mathemagic'' desvi\'{o} hacia interpretaciones matem\'{a}ticas (no ling\"{u}\'{i}sticas)
\item \textbf{La trampa de la criba de Knightmare} --- la criba m\'{a}s prometedora del NPC (``BE A WIT THAN BE A FOOL'') est\'{a} en ingl\'{e}s pero el texto plano es alem\'{a}n, llevando a callejones sin salida
\end{enumerate}

\subsection{Dise\~{n}o de CipSoft}

El cifrado parece ser una codificaci\'{o}n ling\"{u}\'{i}stica genuina en lugar de un rompecabezas puramente matem\'{a}tico:

\begin{itemize}
\item El texto contiene una narrativa coherente con temas consistentes
\item Los nombres propios siguen las convenciones de anagramas documentadas de CipSoft
\item El idioma (alto alem\'{a}n medio/alem\'{a}n) es consistente con que CipSoft sea una empresa alemana
\item El g\'{e}nero ``LEICH'' (eleg\'{i}a/poema) coincide con las tradiciones literarias germ\'{a}nicas medievales
\end{itemize}

La pista del Wrinkled Bonelord sobre ``mathemagic'' puede referirse a las matem\'{a}ticas de la sustituci\'{o}n homof\'{o}nica en lugar de indicar contenido no ling\"{u}\'{i}stico.

\subsection{Limitaciones}

\begin{enumerate}
\item \textbf{Brecha del 5.6\% a nivel de palabra}: \~{}310 caracteres permanecen como bloques distorsionados a nivel de palabra, aunque todos est\'{a}n correctamente decodificados a nivel de letra
\item \textbf{Incertidumbre en la transcripci\'{o}n}: Los libros fueron transcritos del texto dentro del juego por miembros de la comunidad; errores en incluso un solo d\'{i}gito afectan la decodificaci\'{o}n posterior
\item \textbf{Ambig\"{u}edad de nombres propios}: Varios nombres propios (THENAEUT, LGTNELGZ, WRLGTNELNR) resisten la resoluci\'{o}n
\item \textbf{Dificultad de traducci\'{o}n}: El alto alem\'{a}n medio es arcaico y el texto decodificado carece de espacios, haciendo desafiante la interpretaci\'{o}n sem\'{a}ntica
\end{enumerate}

\subsection{Comparaci\'{o}n con Esfuerzos Previos de la Comunidad}

\begin{table}[h]
\centering
\begin{tabular}{ll}
\toprule
Enfoque & Resultado \\
\midrule
Mapeo de coordenadas (TibiaSecrets) & Sin salida coherente \\
Alineamiento de secuencias de ADN (s2ward/469) & Sin convergencia \\
Sustituci\'{o}n en ingl\'{e}s (varios) & Texto sin sentido \\
Codificaci\'{o}n de longitud variable (varios) & Sin mapeo consistente \\
Codificaci\'{o}n de notas musicales (comunidad) & Sin patr\'{o}n \\
\textbf{Nuestro enfoque (homof\'{o}nico + alem\'{a}n + cribas)} & \textbf{94.6\% descifrado} \\
\bottomrule
\end{tabular}
\end{table}

\section{Trabajo Futuro}

\subsection{Criptoan\'{a}lisis Restante}

\begin{itemize}
\item Resolver los \~{}310 caracteres de bloques distorsionados a nivel de palabra
\item Intentar resolver nombres propios sin soluci\'{o}n (THENAEUT, LGTNELGZ, WRLGTNELNR)
\item Decodificar muestras de habla de NPC (Evil Eye, Elder Bonelord) usando el mapeo v7
\item Decodificar el poema 469 de Avar Tar
\end{itemize}

\subsection{Verificaci\'{o}n Dentro del Juego}

\begin{itemize}
\item Verificar las 70 transcripciones de libros contra el texto actual del juego
\item Probar palabras clave decodificadas con NPCs (especialmente el Wrinkled Bonelord)
\item Buscar textos 469 fuera de Hellgate (Isle of Kings, Ferumbras Citadel, Paradox Tower)
\item Investigar si ORANGENSTRASSE corresponde a una ubicaci\'{o}n f\'{i}sica dentro del juego
\end{itemize}

\subsection{An\'{a}lisis Narrativo}

\begin{itemize}
\item Completar la traducci\'{o}n acad\'{e}mica alem\'{a}n $
ightarrow$ ingl\'{e}s del texto completo
\item Identificar el g\'{e}nero literario y potenciales influencias de fuentes en alto alem\'{a}n medio
\item Hacer referencias cruzadas entre la narrativa decodificada y el lore del juego Tibia
\item Determinar si el texto contiene pistas de misiones o mec\'{a}nicas ocultas del juego
\end{itemize}

\subsection{Participaci\'{o}n Comunitaria}

\begin{itemize}
\item Publicar el mapeo y pipeline para verificaci\'{o}n independiente
\item Coordinar con el repositorio s2ward/469
\item Compartir hallazgos con TibiaSecrets y Tales of Tibia
\end{itemize}

\section{Conclusi\'{o}n}

Despu\'{e}s de 31 sesiones de criptoan\'{a}lisis sistem\'{a}tico, hemos logrado un descifrado del 94.6\% a nivel de palabra del cifrado Tibia Bonelord 469 --- la primera soluci\'{o}n conocida en los m\'{a}s de 25 a\~{n}os de historia del cifrado. El texto es un cifrado de sustituci\'{o}n homof\'{o}nica (homophonic substitution) que codifica texto alem\'{a}n usando 98 c\'{o}digos de dos d\'{i}gitos mapeados a 22 letras. El texto descifrado revela una inscripci\'{o}n funeraria bonelord que presenta al Rey Salzberg, los Siervos de Dios, ruinas de piedra antiguas, demonios del bosque y magia r\'{u}nica.

Las innovaciones metodol\'{o}gicas clave fueron: (1) identificar el idioma del texto plano como alem\'{a}n en lugar de ingl\'{e}s; (2) asignaci\'{o}n progresiva de c\'{o}digos por niveles con validaci\'{o}n de cribas; (3) partici\'{o}n de palabras por bolsa de letras con tolerancia de intercambio I$eftrightarrow$E/L; (4) optimizaci\'{o}n de digit-split consciente de concatenaci\'{o}n para libros de longitud impar; y (5) correcciones post-resoluci\'{o}n para artefactos creados por la aplicaci\'{o}n secuencial de anagramas.

A nivel de letra, el texto est\'{a} 100\% descifrado. La brecha restante del 5.4\% a nivel de palabra consiste en letras alemanas correctamente decodificadas en l\'{i}mites de palabras que resisten la segmentaci\'{o}n --- un artefacto del algoritmo de PD, no una brecha criptogr\'{a}fica. El cifrado est\'{a} resuelto.

\appendix

\section{Estructura del Repositorio}

\begin{lstlisting}[style=plain]
tibia-research/
  data/
    books.json              # Source: 70 books as digit strings
    mapping_v7.json         # The mapping (98 codes -> 22 letters)
  scripts/
    core/
      narrative_v3_clean.py # Complete decryption pipeline
    analysis/               # Per-session analysis scripts
  docs/
    narrative_translation.md  # Full translated text (EN/ES)
    roadmap_ingame.md        # In-game verification roadmap
    npc-research.md          # NPC dialogue research
  FINDINGS.md               # Complete 31-session research log
\end{lstlisting}

\section{Reproducci\'{o}n de Resultados}

\begin{lstlisting}[style=bash]
# Decode all 70 books with 94.6% word coverage
python scripts/core/narrative_v3_clean.py
\end{lstlisting}

\section*{References}

\textit{This research was conducted independently using publicly available data from the Tibia community. CipSoft GmbH owns all intellectual property related to Tibia and its content.}

\end{document}
